% Options for packages loaded elsewhere
\PassOptionsToPackage{unicode}{hyperref}
\PassOptionsToPackage{hyphens}{url}
\PassOptionsToPackage{dvipsnames,svgnames,x11names}{xcolor}
\documentclass[
  12pt,
]{article}
\usepackage{xcolor}
\usepackage[letterpaper,margin=1in]{geometry}
\usepackage{amsmath,amssymb}
\setcounter{secnumdepth}{-\maxdimen} % remove section numbering
\usepackage{iftex}
\ifPDFTeX
  \usepackage[T1]{fontenc}
  \usepackage[utf8]{inputenc}
  \usepackage{textcomp} % provide euro and other symbols
\else % if luatex or xetex
  \usepackage{unicode-math} % this also loads fontspec
  \defaultfontfeatures{Scale=MatchLowercase}
  \defaultfontfeatures[\rmfamily]{Ligatures=TeX,Scale=1}
\fi
\usepackage{lmodern}
\ifPDFTeX\else
  % xetex/luatex font selection
\fi
% Use upquote if available, for straight quotes in verbatim environments
\IfFileExists{upquote.sty}{\usepackage{upquote}}{}
\IfFileExists{microtype.sty}{% use microtype if available
  \usepackage[]{microtype}
  \UseMicrotypeSet[protrusion]{basicmath} % disable protrusion for tt fonts
}{}
\usepackage{setspace}
\makeatletter
\@ifundefined{KOMAClassName}{% if non-KOMA class
  \IfFileExists{parskip.sty}{%
    \usepackage{parskip}
  }{% else
    \setlength{\parindent}{0pt}
    \setlength{\parskip}{6pt plus 2pt minus 1pt}}
}{% if KOMA class
  \KOMAoptions{parskip=half}}
\makeatother
\setlength{\emergencystretch}{3em} % prevent overfull lines
\providecommand{\tightlist}{%
  \setlength{\itemsep}{0pt}\setlength{\parskip}{0pt}}
\usepackage{bookmark}
\IfFileExists{xurl.sty}{\usepackage{xurl}}{} % add URL line breaks if available
\urlstyle{same}
\hypersetup{
  colorlinks=true,
  linkcolor={black},
  filecolor={Maroon},
  citecolor={Blue},
  urlcolor={blue},
  pdfcreator={LaTeX via pandoc}}

\author{}
\date{}

\begin{document}

\setstretch{1.15}
\section{Protected Legal Thought in the Age of
AI}\label{protected-legal-thought-in-the-age-of-ai}

\subsection{Work Product and the Civil Rights of Legal
Preparation}\label{work-product-and-the-civil-rights-of-legal-preparation}

\textbf{Author Name Redacted for Blind Review}

\begin{center}\rule{0.5\linewidth}{0.5pt}\end{center}

\subsection{Abstract}\label{abstract}

Federal courts have begun dividing on whether AI-assisted legal
preparation is protected work product. In \emph{United States v.
Heppner}, a court held that a defendant\textquotesingle s AI-generated
legal analysis was discoverable because it was not created under
counsel\textquotesingle s direction. On the same day, in \emph{Warner v.
Gilbarco, Inc.}, a court held that a pro se litigant\textquotesingle s
AI-assisted materials were protected because they reflected her own
litigation strategy. Together, these decisions expose a growing
fracture: materially similar legal preparation receives different
protection based not on the nature of the work, but on the status of the
person performing it.

This Essay argues that this fracture raises an underexamined civil
rights problem. As artificial intelligence becomes a primary tool for
legal preparation, access to protected legal reasoning increasingly
depends on professional credential and platform tier rather than
function. Because effective participation in legal processes---whether
defending against criminal charges, seeking redress, or navigating
administrative systems---requires the ability to prepare without
compelled disclosure, the uneven protection of AI-assisted work risks
entrenching structural inequality across the legal system. Individuals
without access to counsel or enterprise-grade tools may be uniquely
exposed, even as they rely most heavily on AI to exercise legal rights.

This Essay contends that this emerging regime implicates core concerns
of civil rights law, including equal protection, procedural fairness,
and meaningful access to legal institutions. It proposes a functional
rule: work product protection should attach to the nature of legal
preparation in anticipation of proceedings, not the professional
identity of the preparer or the technological channel through which the
work is produced.

\begin{center}\rule{0.5\linewidth}{0.5pt}\end{center}

\subsection{I. Introduction: The Fracture Becomes
Visible}\label{i-introduction-the-fracture-becomes-visible}

On February 10, 2026, two federal courts decided the same question and
reached opposite conclusions.\footnote{\emph{See} United States v.
  Heppner, No. 25-cr-00503-JSR, 2026 WL 436479 (S.D.N.Y. Feb. 17, 2026);
  Warner v. Gilbarco, Inc., No. 2:24-cv-12333, 2026 WL 373043 (E.D.
  Mich. Feb. 10, 2026). Judge Rakoff\textquotesingle s bench ruling was
  issued February 10, 2026; the written opinion followed February 17,
  2026.} The split is not a procedural anomaly. It is the first visible
surface of a structural inequality forming inside AI-mediated legal
systems. The question was whether materials a person created with the
assistance of a publicly available AI platform---legal research, case
analysis, strategic planning---are protected work product, or whether
they are discoverable because the AI platform is not an attorney and its
terms of service disclaim confidentiality. The answer, it turns out,
depends on who the person is.

In \emph{United States v. Heppner}, the U.S. District Court for the
Southern District of New York held that a criminal
defendant\textquotesingle s AI-generated legal documents were not
protected work product.\footnote{Heppner, 2026 WL 436479.} Bradley
Heppner, facing securities fraud charges, had used a consumer AI
platform to create thirty-one documents containing defense strategies,
factual analyses, and legal arguments. He did so without his
attorney\textquotesingle s direction. The court, applying what it called
the "core" of the work-product doctrine---"shelter{[}ing{]} the mental
processes of the attorney, providing a privileged area within which he
can analyze and prepare his client\textquotesingle s case"---found that
materials created by a non-lawyer without attorney involvement fell
outside the doctrine\textquotesingle s shelter.\footnote{\emph{Id.} at
  *3 (quoting United States v. Nobles, 422 U.S. 225, 238 (1975)).} The
AI\textquotesingle s terms of service, which permitted data collection
and disclosure to third parties, reinforced the conclusion that Heppner
had no reasonable expectation of confidentiality. The documents were
deemed fully discoverable.

On the same day, in the Eastern District of Michigan, \emph{Warner v.
Gilbarco, Inc.} reached the opposite result.\footnote{Warner, 2026 WL
  373043.} Sohyon Warner, a pro se plaintiff in an employment
discrimination case, had used consumer AI to prepare her case. When the
defendants moved to compel production of all AI-related materials, the
court denied the motion. Pro se litigants, it reasoned, occupy a unique
dual role---"simultaneously the party and the advocate."\footnote{\emph{Id.}
  at *3.} Their AI-assisted materials reflect their own mental
impressions as the person litigating, and compelling production "would
nullify work-product protection in nearly every modern drafting
environment, a result no court has endorsed."\footnote{\emph{Id.} at *4.}
The AI was a tool, not a person; disclosure to a tool did not waive
protection.

Within weeks, \emph{Morgan v. V2X, Inc.} in the District of Colorado
extended \emph{Warner}\textquotesingle s reasoning, explicitly
distinguishing \emph{Heppner} on the ground that the criminal defendant
there had acted independently of counsel while the pro se civil litigant
was herself the advocate.\footnote{Morgan v. V2X, Inc., No.
  25-cv-01991-SKC-MDB, 2026 WL 864223 (D. Colo. Mar. 30, 2026).} A
doctrinal fracture has opened.

The emerging divide is not about whether AI is reliable, whether its
outputs are accurate, or whether its use constitutes the unauthorized
practice of law. It is about whether the protection of legal
preparation---the ability to research, analyze, and strategize without
surrendering that work to the opposing side---attaches to the nature of
the work or the status of the person performing it. This Essay argues
that the fracture reveals a structural inequality in the distribution of
protected legal thought, one that AI has made visible and that civil
rights law is uniquely positioned to address. At its most extreme, this
fracture leaves crime victims---who are not parties to the criminal
proceeding, not represented by the prosecutor, and not entitled to
appointed counsel---with no work-product framework for their preparation
at all. The core thesis is that work-product protection is increasingly
determined by professional credential and subscription tier rather than
by function, and that this development conditions the exercise of
fundamental rights on status and access.

The argument proceeds in six parts. Part II describes how AI has
transformed access to legal capability without a corresponding expansion
of protection. Part III examines the constitutional infrastructure---the
technical, contractual, and institutional systems---that now mediates
the protection of legal thought. Part IV anchors the fracture in civil
rights law: equal protection, procedural fairness, and access to courts.
Part V proposes a functional rule that would align protection with
preparation rather than status. Part VI identifies materials---including
a recently discovered conservation law---that make operational
verification of protected preparation increasingly feasible. A brief
conclusion returns to the civil rights stakes.

\begin{center}\rule{0.5\linewidth}{0.5pt}\end{center}

\subsection{II. The Transformation of Legal Capability: Capability
Without
Protection}\label{ii-the-transformation-of-legal-capability-capability-without-protection}

Work-product doctrine rests on a factual predicate that AI has
destroyed. The doctrine, recognized in \emph{Hickman v. Taylor} and
codified in Federal Rule of Civil Procedure 26(b)(3), protects
"documents and tangible things that are prepared in anticipation of
litigation or for trial by or for another party or its
representative."\footnote{329 U.S. 495 (1947); Fed. R. Civ. P.
  26(b)(3)(A). The criminal equivalent in New York is CPL § 245.65.} The
historical justification for tying protection to attorney status was
straightforward: only lawyers could generate the kind of legal analysis
the doctrine was designed to shelter. The ability to perform
attorney-quality legal work tracked the professional credential.
Extending protection to the attorney\textquotesingle s
agents---paralegals, investigators, consultants---was a practical
extension, not a theoretical challenge, because those agents operated
under the attorney\textquotesingle s direction. The system was
credential-gated because the capability was credential-gated.

AI has broken that link. A non-lawyer with access to a publicly
available large language model can now research case law, analyze
statutory elements, map evidence to legal standards, identify procedural
requirements, evaluate defense strategies, and draft pleadings at a
level of sophistication that would have required a trained attorney a
decade ago.\footnote{\emph{See generally} Danielle Keats Citron \& Frank
  Pasquale, \emph{The Scored Society: Due Process for Automated
  Predictions}, 89 WASH. L. REV. 1 (2014) (discussing the shift in
  analytical capability from human professionals to automated systems).}
This is not a speculative claim about future technology. It is the
present reality for millions of individuals who navigate the legal
system without counsel---pro se civil litigants, criminal defendants
supplementing appointed counsel\textquotesingle s work, administrative
claimants, and crime victims preparing to engage with
prosecutors.\footnote{The gap in legal representation is well
  documented. \emph{See} Legal Servs. Corp., \emph{The Justice Gap: The
  Unmet Civil Legal Needs of Low-Income Americans} (2022) (finding that
  86\% of civil legal problems reported by low-income Americans received
  inadequate or no legal help).}

The collapse of the credential-to-capability link creates a civil rights
gap. Effective participation in legal proceedings requires preparation,
and preparation under threat of compelled disclosure is not meaningful
preparation. If a person researching their legal rights knows that every
query, every analysis, every drafted strategy can be demanded by the
opposing side---while the opposing side\textquotesingle s lawyer
prepares in the shelter of work-product protection---the playing field
is structurally tilted. The tilt does not arise from differences in the
quality of the work. It arises from the asymmetry of protection.

The \emph{Heppner/Warner} split makes this asymmetry concrete. The
criminal defendant who used AI independently of counsel lost all
protection. The pro se civil plaintiff who used the same kind of AI for
the same kind of preparation retained protection. The decisive variable
was not the nature of the work, the content of the prompts, or the
quality of the output. It was the user\textquotesingle s status relative
to the attorney-client relationship.\footnote{\emph{Compare} Heppner,
  2026 WL 436479, at *4 (no counsel direction, no protection)
  \emph{with} Warner, 2026 WL 373043, at *3 (self-represented party is
  own advocate, protection applies).} For represented parties acting
outside that relationship, AI-assisted work is exposed. For
self-represented parties, AI-assisted work may be sheltered. For crime
victims---who are not parties, not represented by the prosecutor, and
not self-represented litigants---no framework attaches at
all.\footnote{Crime victims are not parties to the criminal proceeding;
  the prosecutor represents the state. \emph{See} Linda R.S. v. Richard
  D., 410 U.S. 614, 619 (1973). Their statutory rights under 18 U.S.C. §
  3771 do not include the right to appointed counsel or to work-product
  protection for their own preparation materials.}

The doctrinal distinction between these categories has a functional
consequence: people who cannot afford counsel or who fall through the
cracks of the existing protection categories are the people most likely
to rely on AI for legal preparation, and their preparation is the least
protected. AI has expanded access to legal capability without expanding
access to protected legal preparation. The capability is distributed;
the protection is not.

\begin{center}\rule{0.5\linewidth}{0.5pt}\end{center}

\subsection{III. Constitutional Infrastructure and the Privatization of
Protected
Thought}\label{iii-constitutional-infrastructure-and-the-privatization-of-protected-thought}

The fracture is compounded by the technological and contractual
architecture through which AI-assisted legal preparation is now
delivered. The infrastructure matters because it is the mechanism
through which protection is distributed.

\subsubsection{A. The Wrapper Ecosystem}\label{a-the-wrapper-ecosystem}

The leading AI tools marketed specifically to the legal profession are
built on the same foundation models available to the public. Harvey, an
AI platform used by major law firms, integrates
Anthropic\textquotesingle s Claude alongside OpenAI\textquotesingle s
GPT models.\footnote{\emph{See} Proskauer Rose LLP, \emph{SDNY Addresses
  Privilege and Work Product Implications of Using Unsecured Public AI
  Tools} (Feb. 18, 2026),
  \url{https://www.proskauer.com/alert/sdny-addresses-privilege-and-work-product-implications-of-using-unsecured-public-ai-tools}
  {[}\url{https://perma.cc/4JWU-BPWQ}{]} (confirming
  Heppner\textquotesingle s use of Anthropic\textquotesingle s Claude);
  Harvey, \emph{Opus 4.6, Now Live in Harvey} (2026),
  \url{https://www.harvey.ai/blog/opus-4-6-now-live-in-harvey}
  {[}\url{https://perma.cc/7JGW-EWA5}{]} (announcing integration of
  Anthropic\textquotesingle s Claude Opus into the Harvey platform);
  \emph{see also} Anthropic, \emph{Customer Story: Harvey},
  \url{https://claude.com/customers/harvey}
  {[}\url{https://perma.cc/E2R7-98JK}{]} (describing
  Harvey\textquotesingle s use of Claude).} CoCounsel, Thomson
Reuters\textquotesingle s AI assistant, is built on
OpenAI\textquotesingle s GPT-4 and holds a contract with the federal
judiciary.\footnote{\emph{See} Press Release, Thomson Reuters,
  \emph{Westlaw Precision with CoCounsel, Practical Law, and CoCounsel
  to be provided to US Federal Courts as the essential information
  provider for the Federal Judiciary} (Apr. 9, 2025),
  \url{https://www.thomsonreuters.com/en/press-releases/2025/april/westlaw-precision-with-cocounsel-practical-law-and-cocounsel-to-be-provided-to-us-federal-courts-as-the-essential-information-provider-for-the-federal-judiciary}
  {[}\url{https://perma.cc/Q5EU-KE4L}{]} (announcing multi-year contract
  with the Administrative Office of the U.S. Courts).} Lexis+ AI and
Protégé draw on multiple foundation models. These enterprise platforms
are sold through subscription tiers that range from tens of thousands to
hundreds of thousands of dollars annually. They are marketed with the
promise that materials created within their environments are protected
as attorney work product.\footnote{\emph{See, e.g.}, LexisNexis,
  \emph{Lexis+ with Protégé} (formerly Lexis+ AI),
  \url{https://www.lexisnexis.com/en-us/products/lexis-plus-ai.page}
  {[}\url{https://perma.cc/J9K6-SWYZ}{]} (describing enterprise-grade
  security and confidentiality for legal work).}

Consumer versions of the same underlying models---Claude, ChatGPT,
Gemini---are available for free or at modest monthly cost. They produce
outputs of materially comparable analytical quality. But the legal
profession\textquotesingle s discourse, reinforced by
\emph{Heppner}\textquotesingle s reasoning, treats materials created on
consumer platforms by non-lawyers as presumptively
unprotected.\footnote{\emph{See} Heppner, 2026 WL 436479, at *4 (citing
  the consumer AI platform\textquotesingle s terms of service as a basis
  for finding no reasonable expectation of confidentiality).}

Publicly available information suggests that the distinction between
enterprise and consumer tools lies primarily in contractual terms, data
handling policies, and integration layers, not in demonstrable
differences in the underlying analytical models.\footnote{\emph{See
  generally} Legaltech Hub, \emph{Generative AI Toolkit},
  \url{https://www.legaltechnologyhub.com/topics/process-improvement/generative-ai-toolkit/}
  {[}\url{https://perma.cc/98ER-ZTT3}{]}; \emph{see also} UC Davis Mabie
  Law Library, \emph{Generative AI Tools and Resources for Law
  Students},
  \url{https://libguides.law.ucdavis.edu/genaiforlawstudents/tools}
  {[}\url{https://perma.cc/GRE9-VP5E}{]} (citing Legaltech
  Hub\textquotesingle s March 2026 mapping of 1,014 generative AI
  product placements in legal tech).} Enterprise agreements typically
prohibit the model provider from training on customer data, include SOC
2 compliance certifications, and add confidentiality provisions. They do
not, however, publicly establish that the model\textquotesingle s
internal computation is segregated between enterprise and consumer data
paths. Publicly available materials do not establish that enterprise and
consumer processing paths are architecturally segregated. The separation
is addressed through data governance promises, not architectural
isolation.\footnote{Publicly available terms of service, privacy
  policies, and industry descriptions reviewed as of April 2026 do not
  identify architectural segregation commitments between enterprise and
  consumer processing paths.}

As a result, the distinction between protected and unprotected
AI-assisted legal preparation appears, in practice and as emerging cases
suggest, to track access tier and user status rather than demonstrable
differences in underlying analytical capability. This is not a
technology distinction dressed in legal clothing. It is a credentialing
product that signals: \emph{this is attorney work, accessed through
attorney channels, therefore the output is protected.} The same model,
accessed through the consumer platform by a non-attorney, produces
output that is not protected---not because the work is different, but
because the access tier identifies the user as a non-attorney.

\subsubsection{B. Power Redistribution}\label{b-power-redistribution}

This architecture directly answers a question the Stanford Law Review
collection poses: how do technological systems redistribute power
between private actors and the state?\footnote{\emph{See} Stanford Law
  Review Online, Call for Submissions: \emph{Technology, Artificial
  Intelligence, and the Future of Civil Rights} (2026),
  \url{https://www.stanfordlawreview.org/online/}
  {[}\url{https://perma.cc/AYZ4-YXSJ}{]}.} The answer, in this context,
is that they redistribute who gets protected legal thought.

Enterprise legal AI platforms have taken publicly available model
capability and wrapped it in a contractual and credentialing layer that
determines whether the constitutional protection of work product
attaches. Courts have not consciously decided to create a two-tier
system. They have, however, applied a doctrine designed for the pre-AI
era in a manner that effectively delegates the protection of legal
reasoning to private subscription architecture. The result is a system
in which protected legal preparation flows through channel partners, not
through constitutional entitlement.

When courts, privilege doctrine, and procedural rules ratify an
arrangement in which protected legal reasoning is functionally available
only through licensed intermediaries and their enterprise tools, the
inequality is state-backed. Private contracts are the mechanism of the
inequality, not a defense to it. The delegation of constitutional
protection into private channel architecture raises questions that
extend beyond privilege doctrine into the domain of civil rights.

\begin{center}\rule{0.5\linewidth}{0.5pt}\end{center}

\subsection{IV. Civil Rights at the Fault Line: Equality, Fairness, and
Access}\label{iv-civil-rights-at-the-fault-line-equality-fairness-and-access}

The AI work-product fracture implicates three foundational civil rights
principles: equal protection, procedural fairness, and access to courts.

\subsubsection{A. Equal Protection and Status-Based
Distinction}\label{a-equal-protection-and-status-based-distinction}

Two persons preparing for litigation---one a lawyer through Harvey, the
other a self-represented litigant through consumer AI---may produce
materially comparable legal analysis. The lawyer\textquotesingle s work
is protected. The self-represented litigant\textquotesingle s work,
under \emph{Heppner}\textquotesingle s logic, may not be. The
distinction does not rest on the function the work serves or the care
with which it was prepared. It rests on the professional status of the
person who prepared it and the platform tier through which it was
produced.

This produces a status-based distinction in the protection of legal
preparation that implicates equal protection principles in ways courts
have not yet confronted directly but cannot indefinitely avoid,
particularly where the distinction systematically disadvantages those
without access to counsel or protected channels of
preparation.\footnote{\emph{Cf.} Griffin v. Illinois, 351 U.S. 12, 19
  (1956) ("There can be no equal justice where the kind of trial a man
  gets depends on the amount of money he has."); \emph{see also} M.L.B.
  v. S.L.J., 519 U.S. 102 (1996) (extending Griffin\textquotesingle s
  principle to civil proceedings involving fundamental interests);
  Douglas v. California, 372 U.S. 353 (1963) (right to counsel on first
  appeal for indigent defendants).} The historical justification for
credential-based asymmetry---that lawyers had unique analytical
capacity---collapses when the same AI model generates materially
comparable analysis for both. When function detaches from credential, a
doctrine that ties protection to credential tracks status, not function.
That classification, when it burdens the population least able to afford
the alternative, is one that civil rights law has reason to scrutinize.

Work-product protection is not merely a procedural convenience. It is a
precondition to the exercise of legal rights. A person who cannot
research, analyze, and strategize in confidence cannot meaningfully
prepare a defense, prosecute a claim, or navigate an administrative
proceeding. In this sense, work-product protection stands in the same
structural relationship to adversarial legal participation that access
to transcripts stands to the right of appeal in \emph{Griffin v.
Illinois}, that access to counsel stands to the right of defense in
\emph{Gideon v. Wainwright}, and that the right of self-representation
stands to the right to present one\textquotesingle s own case in
\emph{Faretta v. California}. Each of those doctrines recognizes that
the right itself is hollow without the conditions that make its exercise
possible. A doctrine that conditions the ability to prepare---to think
legally in private---is functionally part of the right itself, not an
adjunct to it. Griffin\textquotesingle s principle---that "there can be
no equal justice where the kind of trial a man gets depends on the
amount of money he has"---applies directly when what determines that
access is subscription tier rather than financial means per se. The
inequality is the same. Only the mechanism has changed.\footnote{\emph{See}
  Gideon v. Wainwright, 372 U.S. 335 (1963); Faretta v. California, 422
  U.S. 806 (1975); Griffin, 351 U.S. at 19. \emph{See also} Bounds v.
  Smith, 430 U.S. 817 (1977) (prisoners have a fundamental right of
  access to the courts); Lewis v. Casey, 518 U.S. 343 (1996)
  (elaborating on the access-to-courts right and its limitations).}

\subsubsection{B. Procedural Fairness and Technological Due
Process}\label{b-procedural-fairness-and-technological-due-process}

The collection\textquotesingle s framing of procedural fairness is
directly implicated. Preparation exposure is procedural disadvantage.
When one party prepares under shelter---knowing their research,
analysis, and strategy cannot be compelled---and another party prepares
under threat of discovery for performing the identical function, the
proceeding is not procedurally fair. The disadvantage is not in the
quality of the analysis; it is in the safety of the process.\footnote{\emph{See}
  Mathews v. Eldridge, 424 U.S. 319, 335 (1976) (procedural due process
  requires a balancing of the private interest, the risk of erroneous
  deprivation, and the government\textquotesingle s interest).}

This is a form of the technological due process concern that scholars
have identified in other AI-mediated contexts: systems are producing
procedural outcomes that doctrine has not yet recognized as fairness
problems.\footnote{\emph{See} Danielle Keats Citron, \emph{Technological
  Due Process}, 85 WASH. U. L. REV. 1249 (2008); Cary Coglianese \&
  David Lehr, \emph{Regulating by Robot: Administrative Decision Making
  in the Machine-Learning Era}, 105 GEO. L.J. 1147 (2017).} The AI
work-product fracture fits that pattern. The infrastructure through
which legal preparation is conducted is distributing protection
unevenly, and courts have not yet confronted the procedural unfairness
that results.

\subsubsection{C. Access to Courts and the Right to
Petition}\label{c-access-to-courts-and-the-right-to-petition}

Preparation is a prerequisite to participation. The right to petition
the government for redress of grievances,\footnote{U.S. CONST. amend. I.}
the right of self-representation recognized in \emph{Faretta v.
California},\footnote{422 U.S. 806 (1975).} and the due process right to
a fair proceeding\footnote{U.S. CONST. amend. XIV, § 1.} all require the
ability to prepare. If the protection regime penalizes use of the only
affordable preparation tool, access to courts is effectively conditioned
on the ability to afford enterprise-tier legal infrastructure or to
retain counsel.

This burden falls most heavily on populations that are already
structurally underserved: indigent criminal defendants using AI before
counsel is appointed or to supplement appointed
counsel\textquotesingle s work;\footnote{\emph{See} Clio, \emph{Does
  Using AI Waive Attorney-Client Privilege? Breaking Down the United
  States v. Heppner Case}, CLIO BLOG (2026),
  \url{https://www.clio.com/blog/ai-attorney-client-privilege-heppner/}
  {[}\url{https://perma.cc/GE4K-94G4}{]} (analyzing privilege risks
  where defendants use consumer AI without counsel direction).} pro se
civil litigants in housing, employment, and benefits cases;
administrative claimants navigating immigration, Social Security
disability, and licensing proceedings; and crime victims who are not
parties to the criminal case, not represented by the prosecutor, and not
entitled to appointed counsel.\footnote{\emph{See supra} note 12.} For
each, AI is the primary tool for legal preparation. For each, the
protection of that preparation is either absent or ambiguous under
current doctrine.

The fracture extends into every domain where individuals must prepare to
exercise legal rights without the shelter of counsel. Civil rights law
has addressed analogous structural asymmetries---discriminatory hiring
tools, automated benefits denials, proxy discrimination in lending---by
identifying the ways technically neutral rules produce systematically
unequal outcomes.\footnote{\emph{See generally} Solon Barocas \& Andrew
  D. Selbst, \emph{Big Data\textquotesingle s Disparate Impact}, 104
  CALIF. L. REV. 671 (2016) (arguing that facially neutral data
  processing can produce discriminatory outcomes); Anya E.R. Prince \&
  Daniel Schwarcz, \emph{Proxy Discrimination in the Age of Artificial
  Intelligence and Big Data}, 105 IOWA L. REV. 1257 (2020).} The AI
work-product fracture is the same species of problem.

\begin{center}\rule{0.5\linewidth}{0.5pt}\end{center}

\subsection{V. A Rule for the Age of AI: Functional
Protection}\label{v-a-rule-for-the-age-of-ai-functional-protection}

The solution is not to abandon work-product protection or to extend it
indiscriminately. It is to align the protection with the function it
serves: sheltering the preparation that enables meaningful participation
in legal proceedings.

The proposed rule is functional. Work-product protection should attach
to the nature of legal preparation in anticipation of litigation or
legal proceedings, not to the professional identity of the preparer or
the technological channel through which the work is produced.
AI-assisted material retains protection when prepared in anticipation of
proceedings, unless the proponent demonstrates actual adversarially
relevant exposure---beyond ordinary computational processing.\footnote{Courts
  have begun drawing distinctions between technical processing and
  legally meaningful disclosure in other digital contexts. \emph{Cf.}
  Carpenter v. United States, 138 S. Ct. 2206 (2018) (by analogy,
  holding that diminished privacy interests in information shared with a
  third party do not categorically defeat constitutional protection);
  \emph{see also} Nichol v. City of Springfield, No. 6:14-cv-01983-AA
  (D. Or. 2015) (extending work-product protection under Rule 26(b)(3)
  to materials prepared by a non-attorney for a party in anticipation of
  litigation).} This approach builds on decisions that predate the
current AI era, including \emph{Shih v. Petal Card, Inc.}, which held
that work-product protection extends to litigation preparation by a
non-attorney corporate representative acting in anticipation of
suit.\footnote{565 F. Supp. 3d 557, 574 (S.D.N.Y. 2021) (extending
  work-product protection to litigation preparation by a non-attorney
  corporate representative acting in anticipation of suit).}

The doctrine that work-product protection incentivizes the
attorney-client relationship serves important values---competence,
ethics regulation, professional accountability. However, that incentive
argument has no force in the circumstances this rule addresses: where
counsel is unavailable or unaffordable, and the alternative is no
preparation at all. The functional rule protects preparation in the gap,
not in place of counsel.

This rule requires courts to draw a distinction that \emph{Heppner}
collapsed. The mere fact that an AI system processes a prompt does not
necessarily constitute disclosure to an adversary or third party in the
traditional sense, particularly where no evidence of retention or
adversarial access is shown.\footnote{\emph{See} Warner, 2026 WL 373043,
  at *4 ("Generative AI programs are tools, not persons.").} The legally
relevant events are onward sharing, human review by platform staff,
government subpoena of stored data, or actual adversarial access---not
the computational inference itself. Treating the potential for future
training or data retention as an immediate waiver of protection is
equivalent to saying that an email provider\textquotesingle s ability to
scan messages for spam waives attorney-client privilege over the scanned
communications. The doctrine does not require that result, and courts
should not adopt it.

The functional rule is consistent with the text of Federal Rule of Civil
Procedure 26(b)(3)(A), which protects materials prepared "by or for
another party or its representative" without credential-limiting
language.\footnote{Fed. R. Civ. P. 26(b)(3)(A). The
  rule\textquotesingle s text covers "a party or its representative,"
  which on its face does not limit the representative to an attorney.
  \emph{See also Shih}, 565 F. Supp. 3d at 574 (extending protection to
  a non-attorney representative); \emph{Nichol}, No. 6:14-cv-01983-AA
  (same for materials prepared by a non-attorney).} \emph{Warner} and
\emph{Morgan} already take this path. Extending it to all persons
preparing in anticipation of litigation---regardless of whether they are
represented, self-represented, or preparing to participate as a crime
victim in a criminal proceeding---is a doctrinal development, not a
statutory amendment.

Three paths are available. Judicial reinterpretation can adopt the
functional rule as \emph{Warner} and \emph{Morgan} have begun to do.
Rulemaking can amend the Federal Rules of Civil Procedure and parallel
state rules to explicitly extend protection to AI-assisted preparation
by self-represented parties and by parties preparing for any legal
proceeding, regardless of platform tier. And legislation can address the
specific gaps---for example, targeted amendments to the Crime
Victims\textquotesingle{} Rights Act\footnote{18 U.S.C. § 3771. The Act
  currently provides victims with the right to be heard, the right to
  notice, and the right to restitution, but does not address the
  protection of their litigation preparation materials.} or model
provisions for state victims\textquotesingle{} rights campaigns---to
ensure that crime victims\textquotesingle{} AI-assisted preparation
materials are not discoverable by the very defendants whose conduct they
are documenting.\footnote{\emph{See} Marsy\textquotesingle s Law for
  All, \emph{Model Amendment Language}, \url{https://www.marsyslaw.us}
  {[}\url{https://perma.cc/4G8H-2WWF}{]}. The model amendment does not
  include provisions for AI-assisted preparation or work-product
  protection.}

Each path is a civil rights intervention. Each recognizes that the
protection of legal preparation is not a professional courtesy but a
prerequisite to the equal exercise of legal rights.

\begin{center}\rule{0.5\linewidth}{0.5pt}\end{center}

\subsection{VI. Materials for Building: The Infrastructure Beneath the
Rule}\label{vi-materials-for-building-the-infrastructure-beneath-the-rule}

A doctrinal rule is necessary but not itself sufficient. To be
operational, a rule that protects AI-assisted legal preparation requires
the capacity to verify that the protection has been preserved---that the
preparation was conducted in confidence and that its protected character
survived the technological pipeline. The infrastructure for that
verification is beginning to emerge.

One material is a recently discovered conservation law describing how
meaning behaves under transformation.\footnote{{[}Author{]}, \emph{A
  Conservation Law for Commitment in Language Under Transformative
  Compression and Recursive Application} (Mar. 19, 2026) (unpublished
  preprint), \url{https://doi.org/10.5281/zenodo.18792459}.} The
Conservation Law of Commitment establishes that under governed
conditions, the commitment kernel of a signal---its obligations,
prohibitions, permissions, and constraints---is conserved; under
ungoverned conditions, it decays in predictable, measurable ways. The
law is substrate-agnostic. It applies to human transformation, AI
transformation, and any hybrid.

Applying that framework to the work-product context---defining the
kernel as "sheltered preparation in anticipation of proceedings" and
implementing governed transformation protocols---would make the
functional rule\textquotesingle s protection demonstrable, not merely
doctrinal. A person preparing for litigation could demonstrate not only
that they intended their work to be confidential but that the
confidentiality was structurally preserved---through cryptographic
lineage, fidelity verification gates, and recursion testing.\footnote{\emph{See}
  {[}Author{]}, \emph{supra} note 36, at §§ 4--5 (describing the
  six-gate verification protocol and recursion testing methodology).}

These are not speculative technologies. They are operational components
of a governance architecture that has been publicly described and
deposited.\footnote{{[}Author{]}, \emph{Experimental Record for A
  Conservation Law for Commitment in Language Under Transformative
  Compression and Recursive Application} (EXP-001 to EXP-007),
  \url{https://doi.org/10.5281/zenodo.19105225}.} They do not require
opening proprietary models or revealing trade secrets. They certify the
result, not the internals. They allow courts, agencies, and parties to
ask the question that \emph{Heppner} and \emph{Warner} each approached
obliquely: did the protection survive the pipeline?

This Essay does not ask courts to adopt any particular technical
architecture. It argues only that verifiable preservation of
confidentiality and legal meaning is becoming possible---and that the
functional rule proposed here is designed to take advantage of that
possibility as it matures.

\begin{center}\rule{0.5\linewidth}{0.5pt}\end{center}

\subsection{VII. Conclusion: Whose Legal Thought Stays
Protected?}\label{vii-conclusion-whose-legal-thought-stays-protected}

The civil rights project has always been about identifying the ways law
formally grants equality while practically permitting inequality to
persist through structural arrangements. The AI work-product fracture is
one of those arrangements. Courts have not consciously decided to create
a two-tier system of protected legal thought. They have done so
inadvertently, by applying a doctrine designed for the pre-AI world to a
world in which the factual predicate of that doctrine---that legal
analytical capability tracks professional credential---has collapsed.

The result is a system in which lawyers prepare under shelter,
enterprise subscribers prepare under shelter, and everyone else prepares
under threat. The threat falls most heavily on the populations that
civil rights law exists to protect.

Artificial intelligence is reshaping how individuals prepare to assert
and defend their rights. Civil rights law is at an inflection
point---courts are narrowing antidiscrimination doctrine, rethinking
protected classes, and grappling with technologically mediated harm in
every domain. The work-product fracture belongs in that conversation. As
AI becomes the primary tool through which unrepresented individuals
access the legal system, civil rights law must confront a new question:
not only who is protected, but whose legal thought is allowed to remain
protected at all.

\begin{center}\rule{0.5\linewidth}{0.5pt}\end{center}

\subsection{Footnotes}\label{footnotes}

\begin{center}\rule{0.5\linewidth}{0.5pt}\end{center}

\emph{Note: Footnotes 36--38 contain self-citations anonymized for blind
review. Author identifying information provided separately in submission
form fields.}

\begin{center}\rule{0.5\linewidth}{0.5pt}\end{center}

\end{document}
